\section{Schwarz maps on direct sums of matrix algebras}

\begin{definition}[Density states on a finite direct sum]
    \label{def:direct_sum_density_matrices}
    \lean{Matrix.directSumDensityMatrices,
        Matrix.mem_directSumDensityMatrices,
        Matrix.directSumDensityMatrices_isConvex}
    \leanok
    In the existing matrix-family coordinates for
    $\mathcal A=\bigoplus_{k\in I}\MN{d_k}$, its density states are the
    families $A=(A_k)_{k\in I}$ such that every $A_k\succeq0$ and
    $\sum_k\tr(A_k)=1$.
\end{definition}

\begin{theorem}[Convex-extreme states on a finite direct sum]
    \label{thm:extreme_direct_sum_density_matrices}
    \lean{Matrix.pureDirectSumDensityMatrices,
        Matrix.mem_pureDirectSumDensityMatrices,
        Matrix.pureDirectSumDensityMatrices_subset_directSumDensityMatrices,
        Matrix.directSumDensityMatrices_eq_convexHull_pureDirectSumDensityMatrices,
        Matrix.pureDirectSumDensityMatrices_subset_extremePoints,
        Matrix.extremePoints_directSumDensityMatrices,
        Matrix.mem_extremePoints_directSumDensityMatrices_iff}
    \leanok
    \uses{def:direct_sum_density_matrices,
        thm:extreme_density_matrices}
    A density state on $\bigoplus_{k\in I}\MN{d_k}$ is convex-extreme if
    and only if exactly one block is nonzero and that block is a rank-one
    orthogonal projection. These are the relative pure states of the direct
    sum.
\end{theorem}

\begin{proof}\leanok
    Applying the spectral decomposition in each block writes every state as
    a convex combination of one-block rank-one states, with total weight the
    sum of the block traces. For the converse, a convex decomposition of a
    one-block state has zero positive summands in every other block; the
    occupied block is rigid by
    Theorem~\ref{thm:extreme_density_matrices}.
\end{proof}

\begin{definition}[Kraus maps between finite sums of matrix algebras]
    \label{def:kraus_map_between_direct_sums}
    \lean{Matrix.directSumMapExtension, Matrix.IsKrausDirectSumMap}
    \leanok
    Let
    \begin{align}
        \mathcal A
         & =\bigoplus_{i\in I}\MN{n_i},
         &
        \mathcal B
         & =\bigoplus_{j\in J}\MN{m_j}.
        \notag
    \end{align}
    Write $\iota_{\mathcal A},\pi_{\mathcal A}$ and $\iota_{\mathcal B},\pi_{\mathcal B}$ for the corresponding
    block-diagonal embeddings and diagonal-block compressions. The canonical full-matrix extension of $T:\mathcal A\to\mathcal B$ is
    \begin{align}
        \widehat T
         & =\iota_{\mathcal B}\circ T\circ\pi_{\mathcal A}.
        \label{eq:spectral_ds_kraus_extension}
    \end{align}
    The map $T$ is a direct-sum Kraus map when $\widehat T$ has a Kraus representation.
\end{definition}

% This is the finite-sum complete-positivity implication used for the
% transported composites in CPSV16, Appendix C.4, lines 1974--1995.  It does
% not establish Kraus representations for those transported maps; that is the
% tensor-specific preparation and partial-trace step tracked in issue #4438.
\begin{lemma}[Coordinatewise maps, composition, and trace-adjoint Schwarz for direct sums]
    \label{lem:direct_sum_kraus_composition_schwarz}
    \lean{Matrix.isKrausDirectSumMap_piMap,
        Matrix.directSumMapExtension_comp,
        Matrix.IsKrausDirectSumMap.comp, Matrix.isSchwarzDirectSumMap_of_directSumExtension_isSchwarzMap,
        Matrix.IsKrausDirectSumMap.directSumTraceAdjointMap_isSchwarz}
    \leanok
    \uses{def:kraus_map_between_direct_sums}
    A coordinatewise family of completely positive maps between paired
    summands determines a direct-sum Kraus map. For direct-sum Kraus maps
    $T:\mathcal A\to\mathcal B$ and $S:\mathcal B\to\mathcal C$, one has
    $\widehat{S\circ T}=\widehat S\circ\widehat T$, and $S\circ T$ is again a direct-sum Kraus map. If an endomorphism
    $T:\mathcal A\to\mathcal A$ has an extension satisfying
    \begin{align}
        \widehat T(Y^*Y)-\widehat T(Y^*)\widehat T(Y)
         & \succeq0,
        \label{eq:spectral_ds_extension_schwarz}
    \end{align}
    then $T$ satisfies the corresponding inequality in every summand. Finally, if $F:\mathcal A\to\mathcal A$ is a trace-preserving direct-sum
    Kraus map, then its trace adjoint satisfies, for every $X\in\mathcal A$,
    \begin{align}
        F^*(X^*X)-F^*(X^*)F^*(X)
         & \succeq0
        \label{eq:spectral_ds_adjoint_schwarz}
    \end{align}
    in every summand.
\end{lemma}

\begin{proof}\leanok
    \uses{thm:kadison_schwarz}
    Choose a Kraus family $K_{i,r}$ for the map on the $i$-th paired
    summand. The canonical extension of the coordinatewise map is the
    controlled Kraus map
    \begin{align}
        \widehat T(X)
        &=\sum_{i,r}K_{i,r}XK_{i,r}^*.
        \notag
    \end{align}
    Each $K_{i,r}$ is supported on its paired input and output summands, so
    the sum retains precisely the diagonal input blocks.

    Since $\pi_{\mathcal B}\iota_{\mathcal B}=\id_{\mathcal B}$,
    \begin{align}
        \widehat S\,\widehat T
         & =\iota_{\mathcal C}S\pi_{\mathcal B}
        \iota_{\mathcal B}T\pi_{\mathcal A}
        =\iota_{\mathcal C}ST\pi_{\mathcal A}
        =\widehat{S\circ T}.
        \notag
    \end{align}
    If $\widehat T(X)=\sum_iK_iXK_i^*$ and $\widehat S(Y)=\sum_jL_jYL_j^*$, then
    \begin{align}
        \widehat{S\circ T}(X)
         & =\sum_{i,j}(L_jK_i)X(L_jK_i)^*,
        \notag
    \end{align}
    which gives the asserted Kraus representation. For $Y=\iota_{\mathcal A}(X)$, the $i$-th diagonal block of
    (\ref{eq:spectral_ds_extension_schwarz}) is
    \begin{align}
        \left[\widehat T(Y^*Y)
            -\widehat T(Y^*)\widehat T(Y)\right]_i
         & =(T(X^*X))_i-(T(X^*))_i(T(X))_i
        \succeq0,
        \notag
    \end{align}
    which proves the descent assertion. If $F$ preserves the total trace, then $\widehat F$ preserves the ordinary trace, so $(\widehat F)^*$ is
    unital. The adjoint $(\widehat F)^*$ is again a Kraus map, so the
    Kraus-form Kadison--Schwarz inequality gives
    \begin{align}
        (\widehat F)^*(Y^*Y)
        -(\widehat F)^*(Y^*)(\widehat F)^*(Y)
         & \succeq0.
        \notag
    \end{align}
    For $Y=\iota_{\mathcal A}(X)$, use $(\widehat F)^*=\widehat{F^*}$ and project this inequality onto the
    $i$-th diagonal summand. The resulting block identity is
    \begin{align}
        \left[\widehat{F^*}(Y^*Y)
            -\widehat{F^*}(Y^*)\widehat{F^*}(Y)\right]_i
         & =(F^*(X^*X))_i-(F^*(X^*))_i(F^*(X))_i
        \succeq0,
        \notag
    \end{align}
    which is (\ref{eq:spectral_ds_adjoint_schwarz}).
\end{proof}

\begin{theorem}[Mutually inverse CPTP maps give a star equivalence]%
    \label{thm:direct_sum_inverse_cptp_star_equiv}
    \lean{IsKrausCP.kadison_schwarz_of_map_one_eq_one}
    \lean{Matrix.IsSchwarzBetweenDirectSums}
    \lean{Matrix.IsKrausDirectSumMap.isSchwarzBetweenDirectSums}
    \lean{Matrix.IsKrausDirectSumMap.map_conjTranspose_between}
    \lean{Matrix.schwarz_equality_of_mutual_inverse_kraus_direct_sum_maps}
    \lean{Matrix.map_mul_of_mutual_inverse_kraus_direct_sum_maps}
    \lean{Matrix.starAlgEquivOfMutualInverseKrausDirectSumMaps}
    \lean{Matrix.starAlgEquivOfMutualInverseKrausDirectSumMaps_apply}
    \lean{Matrix.starAlgEquivOfMutualInverseKrausDirectSumMaps_symm_apply}
    \lean{Matrix.directSumTraceAdjointMapBetween_comp_eq_id_of_comp_eq_id}
    \lean{Matrix.directSumTraceAdjointMapBetween_map_mul_of_mutual_inverse}
    \lean{Matrix.directSumTraceAdjointStarAlgEquiv}
    \lean{Matrix.directSumTraceAdjointStarAlgEquiv_apply}
    \lean{Matrix.directSumTraceAdjointStarAlgEquiv_symm_apply}
    \leanok
    \uses{def:kraus_map_between_direct_sums,
        lem:direct_sum_rectangular_trace_adjoint} Let $T:\mathcal A\to\mathcal B$ and $S:\mathcal B\to\mathcal A$ be
    mutually inverse completely positive trace-preserving maps between finite direct sums of full matrix algebras. Their trace adjoints are mutually
    inverse unital completely positive maps. Moreover, for all $X,Y\in\mathcal B$,
    \begin{align}
        T^*(XY)
         & =T^*(X)T^*(Y),
         &
        T^*(X^*)
         & =T^*(X)^*,
        \label{eq:spectral_ds_adjoint_multiplicative}
    \end{align}
    and likewise for $S^*$. Thus $T^*$ and $S^*$ determine mutually inverse star-algebra equivalences. This theorem does not yet identify the simple
    summands or their dimensions.
\end{theorem}

\begin{proof}\leanok
    Complete positivity and trace preservation make both trace adjoints unital Schwarz maps. For $F=T^*$ and $G=S^*$, set
    $\Delta_F(X)=F(X^*X)-F(X)^*F(X)\succeq0$. Positivity of $G$ gives $G(\Delta_F(X))\succeq0$. The Schwarz inequality
    for $G$, together with $GF=\id$, gives $-G(\Delta_F(X))\succeq0$. Hence $G(\Delta_F(X))=0$, and injectivity of
    $G$ implies $\Delta_F(X)=0$. Polarization of this equality gives the
    first identity in (\ref{eq:spectral_ds_adjoint_multiplicative}).
    Positivity gives preservation of the adjoint, so the mutually inverse
    maps are star-algebra equivalences. This supplies an algebraic route to
    the block-classification step for which
    \cite[Appendix~C.4, line~1997]{Cirac2016MPDO_arXiv} cites
    \cite[Theorem~8]{WolfPerezGarcia2010Inverse}. The latter instead uses
    extreme states and the positive inverse to identify the blocks, then the
    Schwarz condition to exclude transposition.
\end{proof}

\begin{definition}[Block ideal]\label{def:block_ideal}
    \lean{TwoSidedIdeal.blockIdeal}
    \leanok
    Let $\{R_i\}_{i\in I}$ be a finite family of simple rings. The $i$-th
    \emph{block ideal} in $\prod_{k\in I}R_k$ is
    \begin{align}
        I_i
        &:=\left\{x \in \prod_{k\in I}R_k :
            x_k = 0 \text{ for } k \neq i\right\}.
        \notag
    \end{align}
\end{definition}

\begin{theorem}[Summand matching and dimension preservation]\label{thm:dim_preserved}
    \lean{Matrix.exists_blockEquiv_dim_eq_of_starAlgEquiv_pi_matrix}
    \leanok
    \uses{def:block_ideal}
    Let $\{D_i\}_{i\in I}$ and $\{E_j\}_{j\in J}$ be positive integers.
    A star-algebra isomorphism
    \begin{align}
        \prod_{i\in I} M_{D_i}(\C)
         & \simeq \prod_{j\in J} M_{E_j}(\C)
        \notag
    \end{align}
    determines an equivalence $\sigma:I\simeq J$ such that the isomorphism
    carries the $i$-th block ideal onto the $\sigma(i)$-th block ideal and
    $D_i=E_{\sigma(i)}$ for every $i\in I$. For an automorphism of one product,
    the induced permutation may exchange equal-dimensional block ideals; it
    preserves the dimension along each matched pair.
\end{theorem}

\begin{proof}\leanok
    A star-algebra isomorphism permutes the minimal nonzero central
    idempotents, inducing an equivalence $\sigma:I\simeq J$ between the
    summands. Its restriction to matched ideals is a complex-linear bijection
    between full matrix algebras. Comparing complex dimensions gives
    $D_i^2=E_{\sigma(i)}^2$, and positivity of the block dimensions gives
    $D_i=E_{\sigma(i)}$.
\end{proof}

\begin{theorem}[Paired component maps are bijections]
    \label{thm:component_map_bijective}
    \lean{TwoSidedIdeal.mem_blockIdeal_iff}
    \lean{TwoSidedIdeal.pi_single_mem_blockIdeal}
    \lean{TwoSidedIdeal.blockComponentMap}
    \lean{TwoSidedIdeal.ringEquiv_maps_single_support_between}
    \lean{TwoSidedIdeal.ringEquiv_symm_maps_blockIdeal_between}
    \lean{TwoSidedIdeal.blockComponentMap_injective}
    \lean{TwoSidedIdeal.blockComponentMap_surjective}
    \lean{TwoSidedIdeal.blockComponentMap_bijective}
    \leanok
    \uses{def:block_ideal}
    Let $I$ and $J$ be finite, and let $\{R_i\}_{i\in I}$ and
    $\{S_j\}_{j\in J}$ be families of simple rings. Write $I_i$ and $J_j$ for
    the corresponding block ideals. Suppose that
    $T : \prod_{i\in I}R_i \to \prod_{j\in J}S_j$ is a ring isomorphism and
    that an equivalence $\sigma:I\simeq J$ satisfies
    $T(I_i)=J_{\sigma(i)}$ for every $i\in I$. Then, for each $i\in I$, the map
    \begin{align}
        x
        &\longmapsto
        (T(0,\ldots,0,x,0,\ldots,0))_{\sigma(i)}
        \notag
    \end{align}
    from $R_i$ to $S_{\sigma(i)}$ is bijective.
\end{theorem}

\begin{proof}\leanok
    The matching hypothesis forces elements supported in the $i$-th block to
    remain supported in the $\sigma(i)$-th block. Injectivity follows from
    injectivity of $T$. The forward block equalities and bijectivity of $T$
    give the corresponding equalities for $T^{-1}$, which yield surjectivity.
\end{proof}

\begin{theorem}[Simple-summand matching for inverse CPTP maps]%
    \label{thm:direct_sum_inverse_cptp_block_matching}
    \lean{Matrix.exists_blockEquiv_dim_eq_of_mutual_inverse_kraus_direct_sum_maps}
    \leanok
    \uses{thm:direct_sum_inverse_cptp_star_equiv, thm:dim_preserved}
    Let $\{D_i\}_{i\in I}$ and $\{E_j\}_{j\in J}$ be positive integers, and
    set
    \begin{align}
        \mathcal A
         & =\prod_{i\in I}\MN{D_i},
         &
        \mathcal B
         & =\prod_{j\in J}\MN{E_j}.
        \notag
    \end{align}
    If $T:\mathcal A\to\mathcal B$ and
    $S:\mathcal B\to\mathcal A$ are mutually inverse completely positive
    trace-preserving maps, let $\Phi:\mathcal A\simeq\mathcal B$ be the
    star-algebra isomorphism obtained from their trace adjoints. Write $I_i$
    and $J_j$ for the block ideals of $\mathcal A$ and $\mathcal B$,
    respectively. Then there is an equivalence $\sigma:I\simeq J$ such that,
    for every $i\in I$,
    \begin{align}
        \Phi(I_i)
         & =J_{\sigma(i)},
         &
        D_i
         & =E_{\sigma(i)}.
        \label{eq:spectral_ds_block_matching}
    \end{align}
    This is the simple-summand matching conclusion used in
    \cite[Appendix~C.4, line~1997]{Cirac2016MPDO_arXiv}; the unitary action
    within the paired summands is a separate conclusion.
\end{theorem}

\begin{proof}
    \leanok
    \uses{thm:direct_sum_inverse_cptp_star_equiv, thm:dim_preserved}
    Theorem~\ref{thm:direct_sum_inverse_cptp_star_equiv} gives the
    star-algebra isomorphism $\Phi:\mathcal A\simeq\mathcal B$. Applying
    Theorem~\ref{thm:dim_preserved} to $\Phi$ gives the equivalence $\sigma$,
    the matched block ideals, and the dimension equalities in
    (\ref{eq:spectral_ds_block_matching}).
\end{proof}

The abstract Skolem--Noether theorem below is a generic matrix-algebra fact,
with no tensor-network content; it is relocated here from the single-block
Fundamental Theorem chapter, whose conjugation argument for equal
matrix-product vectors cites it across the chapter boundary.

\begin{theorem}[Skolem--Noether]\label{thm:skolem_noether}
    \lean{Matrix.skolemNoether_matrix}
    \leanok
    Every $\C$-algebra automorphism $f$ of $\MN{D}$ is inner: there is
    $X \in \GL_D(\C)$ with $f(M) = X M X^{-1}$ for all~$M$.
\end{theorem}

\begin{proof}
    \leanok
    Under the identification $\MN{D} \cong \End_{\C}(\C^D)$, every
    algebra automorphism of $\End_{\C}(\C^D)$ is conjugation by an
    invertible linear map. Translating back gives the invertible~$X$.
\end{proof}

\begin{theorem}[Unitary Skolem--Noether theorem for complex matrices]%
    \label{thm:unitary_skolem_noether_matrix}
    \lean{Matrix.exists_unitary_conj_of_starAlgEquiv_matrix}
    \leanok
    Every star-algebra automorphism $\Phi$ of a nonzero full complex matrix
    algebra is implemented by a unitary: there is a unitary matrix $U$ such
    that $\Phi(X)=UXU^*$ for every matrix $X$.
\end{theorem}

\begin{proof}\leanok
    \uses{thm:skolem_noether}
    By Skolem--Noether, there is an invertible matrix $P$ such that
    $\Phi(X)=PXP^{-1}$. Preservation of the adjoint implies that $P^*P$
    commutes with every matrix, hence $P^*P=c\Id$ for a scalar $c$.
    Positivity and invertibility give $c>0$. Therefore $U=c^{-1/2}P$ is
    unitary and implements the same automorphism.
\end{proof}

\begin{theorem}[Unitary action within matched simple summands]%
    \label{thm:direct_sum_star_equiv_block_unitary}
    \lean{Matrix.blockComponentStarAlgEquiv}
    \lean{Matrix.matrixReindexStarAlgEquiv}
    \lean{Matrix.exists_unitary_block_implementers_of_starAlgEquiv_pi_matrix}
    \lean{Matrix.exists_blockEquiv_dim_eq_unitary_of_starAlgEquiv_pi_matrix}
    \lean{Matrix.exists_blockEquiv_dim_eq_unitary_of_mutual_inverse_kraus_direct_sum_maps}
    \leanok
    Retain an equivalence $\sigma:I\simeq J$ which matches the block ideals
    of a star-algebra isomorphism
    $\Phi:\prod_{i\in I}\MN{D_i}\simeq\prod_{j\in J}\MN{E_j}$,
    together with equalities $D_i=E_{\sigma(i)}$. Then there are unitaries
    $U_i\in\MN{E_{\sigma(i)}}$ such that, for every
    $X\in\MN{D_i}$,
    \begin{align}
        \Phi(0,\ldots,0,X,0,\ldots,0)_{\sigma(i)}
         & =U_i\,\iota_i(X)\,U_i^*,
        \label{eq:spectral_ds_block_unitary}
    \end{align}
    where $\iota_i$ is the reindexing induced by the retained dimension
    equality. In particular, this conclusion applies to the trace-adjoint
    star-algebra isomorphism of any mutually inverse completely positive
    trace-preserving pair. This is the blockwise-unitary statement in
    \cite[Appendix~C.4, line~1997]{Cirac2016MPDO_arXiv}; it does not assert
    the later MPDO multiplicity or coefficient relations.
\end{theorem}

\begin{proof}\leanok
    \uses{thm:unitary_skolem_noether_matrix, thm:dim_preserved,
        thm:direct_sum_inverse_cptp_star_equiv}
    The restriction of $\Phi$ to a matched block is a star-algebra
    isomorphism. After reindexing by $D_i=E_{\sigma(i)}$, it is an
    automorphism of a full complex matrix algebra. Skolem--Noether writes
    this automorphism as conjugation by an invertible matrix $P$.
    Preservation of the adjoint implies that $P^*P$ commutes with every
    matrix and is therefore a positive scalar multiple of the identity.
    Dividing $P$ by the positive square root of this scalar produces the
    unitary in (\ref{eq:spectral_ds_block_unitary}).
\end{proof}

\begin{lemma}[Total trace under unitary action on matched summands]
    \label{lem:direct_sum_block_unitary_trace_preserving}
    \lean{Matrix.isTracePreservingBetweenDirectSums_of_unitary_block_implementers}
    \leanok
    Let $\Phi:\mathcal A\simeq\mathcal B$ match the simple summands by an
    equivalence $\sigma:I\simeq J$, with $D_i=E_{\sigma(i)}$. Suppose that
    there are unitaries $U_i$ satisfying
    $\Phi(0,\ldots,0,X,0,\ldots,0)_{\sigma(i)}
        =U_i\iota_i(X)U_i^*$.
    Then $\Phi$ preserves the total block trace.
\end{lemma}

\begin{proof}\leanok
    The image of a matrix supported on the $i$-th summand is supported on the
    $\sigma(i)$-th summand. Unitary conjugation and reindexing preserve the
    trace, so
    \begin{align}
        \sum_j\tr
        (\Phi(0,\ldots,0,X,0,\ldots,0)_j)
         & =\tr(X).
        \notag
    \end{align}
    Summing this identity over the source summands proves the result.
\end{proof}

\begin{theorem}[Unitary action of an invertible CPTP map]
    \label{thm:direct_sum_inverse_cptp_forward_unitary}
    \lean{Matrix.%
        exists_blockEquiv_dim_eq_unitary_forward_of_mutual_inverse_kraus_direct_sum_maps}
    \leanok
    Let $T:\mathcal A\to\mathcal B$ and $S:\mathcal B\to\mathcal A$ be
    mutually inverse completely positive trace-preserving maps between finite
    products of nonzero full matrix algebras. There are an equivalence
    $\sigma:I\simeq J$, equalities $D_i=E_{\sigma(i)}$, and unitaries
    $U_i\in\MN{E_{\sigma(i)}}$ such that
    \begin{align}
        T(0,\ldots,0,X,0,\ldots,0)
         & =(0,\ldots,0,U_i\iota_i(X)U_i^*,0,\ldots,0),
        \label{eq:spectral_ds_forward_unitary}               \\
        S(0,\ldots,0,Y,0,\ldots,0)
         & =(0,\ldots,0,\iota_i^{-1}(U_i^*YU_i),0,\ldots,0),
        \label{eq:spectral_ds_inverse_unitary}
    \end{align}
    with the nonzero entry in
    (\ref{eq:spectral_ds_forward_unitary}) in position $\sigma(i)$ and the
    entries on the left and right of
    (\ref{eq:spectral_ds_inverse_unitary}) in positions $\sigma(i)$ and $i$,
    respectively. These identities hold for every $i\in I$,
    $X\in\MN{D_i}$, and $Y\in\MN{E_{\sigma(i)}}$.
\end{theorem}

\begin{proof}\leanok
    \uses{thm:direct_sum_inverse_cptp_star_equiv,
        thm:direct_sum_star_equiv_block_unitary,
        lem:direct_sum_block_unitary_trace_preserving,
        lem:direct_sum_trace_adjoint_star_equiv_inverse,
        lem:direct_sum_rectangular_trace_adjoint}
    Let $\Phi=S^*:\mathcal A\simeq\mathcal B$ be the star-algebra
    isomorphism obtained from the inverse pair.
    Theorem~\ref{thm:direct_sum_star_equiv_block_unitary} supplies $\sigma$,
    the dimension equalities, and the unitaries for $\Phi$.
    Lemma~\ref{lem:direct_sum_block_unitary_trace_preserving} shows that
    $\Phi$ preserves the total trace, hence
    Lemma~\ref{lem:direct_sum_trace_adjoint_star_equiv_inverse} gives
    $\Phi^*=\Phi^{-1}$. On the other hand, involutivity of the trace adjoint
    gives $\Phi^*=S$. Therefore $S=\Phi^{-1}$, and the two inverse laws imply
    $T=\Phi$. The unitary formula for $\Phi$ gives
    (\ref{eq:spectral_ds_forward_unitary}), while the inverse formula for
    $S=\Phi^{-1}$ gives (\ref{eq:spectral_ds_inverse_unitary}) with the same
    summand matching and the same unitaries.
\end{proof}

% Scope restriction documented in
% docs/paper-gaps/wolf_theorem6_14_fixed_point_projection_gap.tex.
\begin{theorem}[Weighted block form of the adjoint Ces\`aro projection on full support]%
    \label{thm:adjoint_mean_ergodic_projection_weighted_block_form}
    \lean{IsPositiveMap.exists_block_densities_of_adjoint_meanErgodicProjection}
    \leanok
    \uses{thm:positive_unital_adjoint_mean_ergodic_retraction,
        thm:fixedPointsStarSubalgebra,
        thm:positive_retraction_finite_star_algebra}
    Let $T:\MN{D}\to\MN{D}$ be positive and trace preserving, suppose that
    $T^*$ satisfies the Schwarz inequality, and suppose that there is a
    positive definite matrix $\rho$ satisfying $T(\rho)=\rho$. Let $P_T$ be
    the mean-ergodic projection of $T$. There are positive integers
    $d_k,m_k$, a unitary $U$, and density matrices
    $\sigma_k\in\MN{m_k}$ such that
    \begin{align}
        U^*\Fix(T^*)U
         & =\bigoplus_k\Id_{m_k}\otimes\MN{d_k},
        \label{eq:spectral_ds_adjoint_fixed_blocks}                                   \\
        P_T^*(A)
         & =U\left(\bigoplus_k\Id_{m_k}\otimes
             \tr_{m_k}\!\left((\sigma_k\otimes\Id_{d_k})(U^*AU)_{kk}\right)\right)U^*
        \label{eq:spectral_ds_adjoint_projection}
    \end{align}
    for every $A\in\MN{D}$. This is the full-support
    conditional-expectation step in the proof of
    \cite[Theorem~6.14]{Wolf2012Quantum}.
\end{theorem}

\begin{proof}
    \leanok
    \uses{thm:positive_unital_adjoint_mean_ergodic_retraction,
        thm:fixedPointsStarSubalgebra,
        thm:positive_retraction_finite_star_algebra}
    The map $P_T^*$ is a positive unital idempotent map whose range is
    $\Fix(T^*)$. For $A\in\Fix(T^*)$, the Schwarz inequality gives
    $T^*(A^*A)\geq T^*(A)^*T^*(A)=A^*A$, while
    \begin{align}
        \tr\!\left(\rho T^*(A^*A)\right)
         & =\tr\!\left(T(\rho)A^*A\right)
        =\tr(\rho A^*A).
        \notag
    \end{align}
    Since $\rho$ is positive definite, these relations imply
    $T^*(A^*A)=A^*A$. Hence $\Fix(T^*)$ is a $*$-subalgebra, and $P_T^*$ is
    a positive retraction onto it. The blockwise classification of positive
    retractions gives the unitary, the matrix dimensions, the density
    matrices $\sigma_k$, and
    (\ref{eq:spectral_ds_adjoint_fixed_blocks})--%
    (\ref{eq:spectral_ds_adjoint_projection}).
\end{proof}

% Scope restriction and factor-order convention documented in
% docs/paper-gaps/wolf_theorem6_14_fixed_point_projection_gap.tex.
\begin{theorem}[Density-block form of the Ces\`aro projection on full support]%
    \label{thm:mean_ergodic_projection_density_block_form}
    \lean{IsPositiveMap.exists_block_densities_of_meanErgodicProjection}
    \leanok
    \uses{thm:adjoint_mean_ergodic_projection_weighted_block_form}
    Under the hypotheses of
    Theorem~\ref{thm:adjoint_mean_ergodic_projection_weighted_block_form},
    the mean-ergodic projection of $T$ satisfies
    \begin{align}
        P_T(B)
         & =U\left(\bigoplus_k\sigma_k\otimes
             \tr_{m_k}((U^*BU)_{kk})\right)U^*.
        \label{eq:spectral_ds_density_projection}
    \end{align}
    Consequently,
    \begin{align}
        \Fix(T)
         & =U\left(\bigoplus_k\sigma_k\otimes\MN{d_k}\right)U^*.
        \label{eq:spectral_ds_density_fixed_blocks}
    \end{align}
    Here each summand is written on
    $\C^{m_k}\otimes\C^{d_k}$, so $\tr_{m_k}$ traces the first factor. This
    is the full-support part of Equation~(6.63) in
    \cite[Theorem~6.14]{Wolf2012Quantum}; the density matrices are not yet
    asserted to be positive definite.
\end{theorem}

\begin{proof}
    \leanok
    \uses{thm:adjoint_mean_ergodic_projection_weighted_block_form}
    For one summand, cyclicity of the trace and the defining property of the
    partial trace give
    \begin{align}
        \tr\!\left((\sigma_k\otimes\tr_{m_k}B)A\right)
         & =\tr\!\left(B\bigl(\Id_{m_k}\otimes
        \tr_{m_k}((\sigma_k\otimes\Id_{d_k})A)\bigr)\right).
        \notag
    \end{align}
    Summing over the diagonal summands and applying the unitary change of
    basis proves (\ref{eq:spectral_ds_density_projection}) by
    nondegeneracy of the trace pairing. Its range is
    (\ref{eq:spectral_ds_density_fixed_blocks}) because
    $\tr_{m_k}(\sigma_k\otimes X)=X$.
\end{proof}

% Scope restriction documented in
% docs/paper-gaps/wolf_theorem6_14_fixed_point_projection_gap.tex.
\begin{theorem}[Density-block form of fixed points on a finite direct sum]%
    \label{thm:direct_sum_fixed_points_density_block_form}
    \lean{Matrix.IsPositiveDirectSumMap.exists_block_densities_of_fixedPoints}
    \leanok
    \uses{thm:direct_sum_full_matrix_extension_compatibility,
        thm:mean_ergodic_projection_density_block_form}
    Let $\mathcal A=\bigoplus_{i\in I}\MN{n_i}$ be a finite direct sum.
    Suppose that $T:\mathcal A\to\mathcal A$ is positive and preserves the
    total trace, that its trace adjoint satisfies the Schwarz inequality,
    and that $T$ fixes a family $\rho=(\rho_i)_i$ with every $\rho_i$
    positive definite. Then there are positive integers $d_k,m_k$, density
    matrices $\sigma_k\in\MN{m_k}$, a unitary $U$ on
    $\bigoplus_i\C^{n_i}$, and a reindexing of this space by
    $\bigoplus_k(\C^{m_k}\otimes\C^{d_k})$ such that $T(A)=A$ if and only
    if there are matrices $X_k\in\MN{d_k}$ satisfying
    \begin{align}
        U^*\iota(A)U
         & =\bigoplus_k\sigma_k\otimes X_k.
        \label{eq:spectral_ds_fixed_point_form}
    \end{align}
    Here $\iota$ is the block-diagonal embedding. This is the
    finite-direct-sum fixed-point step used in
    \cite[Appendix~C.4, lines~1980--1995]{Cirac2016MPDO_arXiv}; it does not
    assert the channel hypotheses for the particular sector maps in that
    argument. This is the full-support restriction: the support reduction
    and complementary zero summand of the general theorem remain open.
\end{theorem}

\begin{proof}
    \leanok
    \uses{thm:direct_sum_full_matrix_extension_compatibility,
        thm:mean_ergodic_projection_density_block_form}
    Extend $T$ by diagonal compression and block-diagonal embedding. The
    canonical extension is positive and trace preserving, its trace adjoint
    satisfies the Schwarz inequality, and the embedded family $\iota(\rho)$
    is a positive-definite fixed point. After choosing a numbered basis of
    the finite-dimensional ambient space, apply the full-support
    density-block theorem. Simultaneous reindexing preserves positivity,
    trace preservation, the Schwarz inequality, the trace adjoint, and
    positive definiteness. Reindexing the resulting unitary and block
    decomposition back to the original space gives
    (\ref{eq:spectral_ds_fixed_point_form}).
\end{proof}

\section{$*$-subalgebras with scalar centre}

The decomposition of Theorem~\ref{thm:starSubalgebra_wedderburnArtin} has a
single factor exactly when the centre of the subalgebra is as small as it can
be. This is the building block of the structure theorem for finite-dimensional
$C^*$-algebras \cite[lines~2082--2098]{SchumacherWerner2004}: cutting an
algebra by the minimal projections of its centre leaves summands whose centre
consists of the multiples of the unit, and each such summand is a full matrix
algebra. The support algebras of a one-dimensional cellular automaton are
identified with full matrix algebras in exactly this
way~\cite[lines~1278--1282]{GrossNesmeVogtsWerner2012}.

The source states the decomposition of an abstract finite-dimensional
$C^*$-algebra into a direct sum of full matrix algebras. What follows treats a
$*$-subalgebra of a full matrix algebra and proves the single-block statement,
not that decomposition; see \cite{gap:sw04_csform_matrix_scope}.

\begin{definition}[Conjugation by an invertible matrix]%
    \label{def:matrix_innerAlgEquiv}
    \lean{Matrix.innerAlgEquiv}
    \lean{Matrix.innerAlgEquiv_apply}
    \leanok
    Let $R,K\in\MN{m}$ satisfy $RK=KR=\Id$. Conjugation
    \begin{align}
        X
         & \longmapsto RXK
        \label{eq:scalar_centre_conjugation}
    \end{align}
    is a $\C$-algebra automorphism of $\MN{m}$, with inverse $X\mapsto KXR$.
\end{definition}

An isomorphism produced by the Wedderburn theory of
Theorem~\ref{thm:starSubalgebra_wedderburnArtin} respects only the ring
structure. The next theorem corrects such an isomorphism, by a conjugation
(\ref{eq:scalar_centre_conjugation}), into one that also intertwines the
adjoints.

\begin{theorem}[Adjoint-preserving correction of an isomorphism onto a matrix
        algebra]%
    \label{thm:starSubalgebra_starAlgEquiv_of_algEquiv_matrix}
    \lean{StarSubalgebra.exists_starAlgEquiv_of_algEquiv_matrix}
    \lean{Matrix.star_dotProduct_conjTranspose_mulVec}
    \lean{Matrix.conjTranspose_vecMulVec_mul_mul_vecMulVec}
    \leanok
    \uses{def:matrix_innerAlgEquiv}
    Let $S$ be a $*$-subalgebra of $\MD$ and let $\varphi:S\to\MN{m}$, with
    $m\geq1$, be an isomorphism of $\C$-algebras. Then $S$ and $\MN{m}$ are
    isomorphic as $*$-algebras.
\end{theorem}

\begin{proof}
    \leanok
    \uses{def:matrix_innerAlgEquiv, thm:skolem_noether}
    Transporting the adjoint of $S$ through $\varphi$ gives the map
    $Z\mapsto\varphi\bigl(\varphi^{-1}(Z)^\dagger\bigr)^\dagger$, a
    $\C$-algebra automorphism of $\MN{m}$. By
    Theorem~\ref{thm:skolem_noether} it is conjugation by an invertible $X$;
    with $H=X^\dagger$ and $K=(X^{-1})^\dagger$, so that $HK=KH=\Id$,
    \begin{align}
        \varphi(y^\dagger)
         & =K\,\varphi(y)^\dagger H,
        \qquad y\in S.
        \label{eq:scalar_centre_transported_adjoint}
    \end{align}
    Applying (\ref{eq:scalar_centre_transported_adjoint}) twice gives
    $Z=KH^\dagger ZK^\dagger H$ for every $Z\in\MN{m}$, so $KH^\dagger$
    commutes with every matrix and is therefore a scalar $c_0$, that is,
    $H^\dagger=c_0H$.

    Fix a coordinate vector $w$, put $W=ww^\dagger$ and
    $c=\tr\varphi^{-1}(KW)$. For $v\neq0$ set $Y=vw^\dagger$ and
    $y=\varphi^{-1}(Y)$. Then $y\neq0$, and
    (\ref{eq:scalar_centre_transported_adjoint}) together with
    \begin{align}
        Y^\dagger HY
         & =(v^\dagger Hv)\,W
        \label{eq:scalar_centre_rank_one_sandwich}
    \end{align}
    gives $y^\dagger y=(v^\dagger Hv)\,\varphi^{-1}(KW)$ and hence
    \begin{align}
        \tr(y^\dagger y)
         & =(v^\dagger Hv)\,c
           >0.
        \label{eq:scalar_centre_positive_pairing}
    \end{align}
    So $c\neq0$, and $v^\dagger Hv\neq0$ for every $v\neq0$. Conjugating the
    Hermitian form of $H$ turns $H$ into $H^\dagger$, so
    $\overline{v^\dagger Hv}=c_0\,v^\dagger Hv$; combined with the reality of
    the left side of (\ref{eq:scalar_centre_positive_pairing}) this gives
    $\overline{c}\,c_0=c$, which says that $cH$ is Hermitian. By
    (\ref{eq:scalar_centre_positive_pairing}) its Hermitian form is positive
    on nonzero vectors, so $cH$ is positive definite.

    Replacing $H$ by $cH$ and $K$ by $c^{-1}K$ preserves
    (\ref{eq:scalar_centre_transported_adjoint}) and the relations
    $HK=KH=\Id$. Write $cH=R^\dagger R$ with $R$ invertible. Then
    $K=R^{-1}(R^{-1})^\dagger$, and substituting into
    (\ref{eq:scalar_centre_transported_adjoint}) gives
    \begin{align}
        R\,\varphi(y^\dagger)\,R^{-1}
         & =\bigl(R\,\varphi(y)\,R^{-1}\bigr)^\dagger.
        \label{eq:scalar_centre_corrected_isomorphism}
    \end{align}
    So the conjugated isomorphism $y\mapsto R\,\varphi(y)\,R^{-1}$ preserves
    adjoints.
\end{proof}

\begin{theorem}[A $*$-subalgebra with scalar centre is simple]%
    \label{thm:starSubalgebra_isSimpleRing_of_isCentral}
    \lean{StarSubalgebra.isSimpleRing_of_isCentral}
    \leanok
    Let $S$ be a $*$-subalgebra of $\MD$ with $D\geq1$, and suppose the centre
    of $S$ consists of the multiples of the unit. Then $S$ is a simple ring.
\end{theorem}

\begin{proof}
    \leanok
    \uses{thm:starSubalgebra_wedderburnArtin}
    By Theorem~\ref{thm:starSubalgebra_wedderburnArtin} there is an
    isomorphism $S\cong\prod_{k=1}^{m}\MN{d_k}$ with every $d_k\geq1$. The
    element that is the unit in the $k$-th coordinate and zero elsewhere is
    central, hence a multiple $\lambda$ of the unit. If $m\geq2$, comparing
    the $k$-th coordinate with another one gives $\lambda=1$ and $\lambda=0$,
    which is impossible. If $m=0$ the product is trivial, whereas $S$ contains
    the unit of $\MD$, which is nonzero because $D\geq1$. So $m=1$ and
    $S\cong\MN{d_1}$ is simple.
\end{proof}

\begin{theorem}[A $*$-subalgebra with scalar centre is a full matrix algebra]%
    \label{thm:starSubalgebra_starAlgEquiv_matrix_of_isCentral}
    \lean{StarSubalgebra.exists_starAlgEquiv_matrix_of_isCentral}
    \leanok
    \uses{thm:starSubalgebra_isSimpleRing_of_isCentral}
    Let $S$ be a $*$-subalgebra of $\MD$ with $D\geq1$, and suppose the centre
    of $S$ consists of the multiples of the unit. Then there is $r\geq1$ and a
    $*$-isomorphism $S\cong\MN{r}$.
\end{theorem}

\begin{proof}
    \leanok
    \uses{thm:starSubalgebra_isSimpleRing_of_isCentral,
        thm:starSubalgebra_starAlgEquiv_of_algEquiv_matrix}
    By Theorem~\ref{thm:starSubalgebra_isSimpleRing_of_isCentral} the ring $S$
    is simple, so the Wedderburn theorem for simple algebras over the
    algebraically closed field $\C$ gives $r\geq1$ and an isomorphism of
    $\C$-algebras $S\cong\MN{r}$. It preserves adjoints after the correction
    of Theorem~\ref{thm:starSubalgebra_starAlgEquiv_of_algEquiv_matrix}.
\end{proof}

\begin{lemma}[Dimension of a subalgebra presented as a full matrix algebra]%
    \label{lem:starSubalgebra_finrank_matrix_presentation}
    \lean{StarSubalgebra.finrank_eq_mul_self_of_starAlgEquiv_matrix}
    \leanok
    Let $S$ be a $*$-subalgebra of $\MD$ admitting a $*$-isomorphism
    $S\cong\MN{r}$. Then $\dim_{\C}S=r^2$.
\end{lemma}

\begin{proof}
    \leanok
    A $*$-isomorphism is in particular a $\C$-linear isomorphism, and
    $\dim_{\C}\MN{r}=r^2$.
\end{proof}
